\chapter{Datasets, Models, and Metrics}
\label{ch:datasets_models_metrics}

% For AI-assisted explanatory figures in this chapter, use \aifigurecaption{...}
% so the standardized disclosure is included directly in the caption.

This work focuses on two pretrained image tokenizers: VQGAN
~\cite{esser2021taming} and LlamaGen VQ-16~\cite{sun2024llamagen}. This chapter
defines these two tokenizers, the evaluation datasets, and the metrics used to
measure reconstruction quality, codebook use, and spatial behavior.

\section{Models Under Study}
\label{sec:models_under_study}

These two tokenizers share the same external task, mapping $256\times256$
images to discrete token grids. They use the same VQGAN-style convolutional
encoder--decoder backbone, while their quantizer projections, codebook
geometry, and training recipes differ~\cite{esser2021taming,sun2024llamagen}.

VQGAN uses an encoder, vector quantizer, and decoder with a convolutional
architecture and attention at the lowest-resolution level
~\cite{esser2021taming}. Section~\ref{sec:vqgan_taming} describes the model in
more detail. The pretrained version used in this study was trained with a
reconstruction, perceptual, adversarial, and codebook objective. It has a
downsampling factor $f=16$, so a $256\times256$ image becomes a $16\times16$
token grid. Its codebook has $K=16{,}384$ entries with dimension $256$.

LlamaGen provides two pretrained versions, VQ-16 and VQ-8, with downsampling
factors of 16 and 8, respectively~\cite{sun2024llamagen,llamagengithub}. To
maintain the same token-grid size as VQGAN, these experiments use VQ-16. It maps
$256\times256$ images to a $16\times16$ grid with a nominal codebook of
$K=16{,}384$. Its codebook embeddings have dimension $d=8$. Before
nearest-neighbor lookup, the projected encoder vectors and codebook vectors are
normalized to unit L2 norm.

Table~\ref{tab:model_config_summary} summarizes the configurations of the two
evaluated tokenizers.

\begin{table}[htbp]
    \centering
    \renewcommand{\arraystretch}{1.2}
    \begin{tabularx}{\textwidth}{>{\raggedright\arraybackslash}p{0.22\textwidth} >{\raggedright\arraybackslash}X >{\raggedright\arraybackslash}X}
        \textbf{Property} & \textbf{VQGAN} & \textbf{LlamaGen} \\ \hline
        Source & Taming Transformers~\cite{esser2021taming} & LlamaGen~\cite{sun2024llamagen} \\
        Main tokenizer & ImageNet first stage, $f=16$ & \texttt{VQ-16} \\
        Input resolution & $256\times256$ & $256\times256$ \\
        Token grid & $16\times16$ & $16\times16$ \\
        Encoder/decoder & VQGAN-style convolutional backbone & VQGAN-style convolutional backbone \\
        Nominal codebook size & $K=16{,}384$ & $K=16{,}384$ \\
        Embedding dimension & $256$ & $8$ \\
        Codebook normalization & none & L2-normalized lookup \\
        Training loss components & pixel, perceptual, adversarial, codebook/commitment & pixel, perceptual, adversarial, codebook/commitment \\
    \end{tabularx}
    \caption{Tokenizer configurations compared in the thesis. Both models produce the same token-grid size and nominal vocabulary size and share the encoder--decoder backbone. LlamaGen uses an L2 pixel loss with LPIPS and PatchGAN terms~\cite{sun2024llamagen}; the objectives differ in details and weights, rather than in the presence of perceptual or adversarial supervision.}
    \label{tab:model_config_summary}
\end{table}

\section{Datasets}
\label{sec:datasets}

The ImageNet-1K validation split is used for reconstruction baselines, codebook
usage, global perturbation robustness, and locality experiments. It belongs to
the training domain of both selected checkpoints and appears in their published
evaluations~\cite{esser2021taming,sun2024llamagen}. ImageNet-V2 is also used in
Chapter~\ref{ch:rq4} for decoder locality. Chapter~\ref{ch:rq5} studies
distribution shifts using ImageNet-V2, ImageNet-Sketch, ObjectNet, BloodMNIST,
OrganAMNIST, and RVL-CDIP.

\subsection{ImageNet}
\label{subsec:imagenet_validation}

ImageNet was introduced as a large-scale visual recognition dataset organized around the WordNet noun hierarchy~\cite{deng2009imagenet}. It was collected to provide broad object-category coverage at a scale that earlier recognition datasets did not offer. The ImageNet Large Scale Visual Recognition Challenge (ILSVRC) later standardized a 1,000-class classification benchmark, commonly called ImageNet-1K~\cite{russakovsky2015imagenet}.

ImageNet is a standard reference dataset for the evaluated VQ tokenizers
~\cite{esser2021taming,sun2024llamagen}. We use the ImageNet-1K validation split
because it matches their training domain. It contains 50,000 images, with 50
images for each of 1,000 object classes. The images are converted to RGB,
resized, center-cropped to $256\times256$, and normalized for each tokenizer.
No test-time data augmentation is used.

\subsection{Distribution-Shifted Datasets}
\label{subsec:shifted_datasets}

These datasets evaluate the tokenizers beyond their training domain. We
distinguish ImageNet-aligned shifts from datasets with different image content.
ImageNet-V2 supports the decoder-locality analysis in Chapter~\ref{ch:rq4}; all
datasets in this section support the reconstruction and codebook analyses in
Chapter~\ref{ch:rq5}.

\paragraph{ImageNet-aligned shifts.}
ImageNet-V2 was collected years after the original ImageNet construction as a
new test set for the same 1,000 classes~\cite{recht2019imagenet}. It keeps the
label space and natural-image setting close to ImageNet while changing the
sampled images, so it represents a near-domain shift.

ImageNet-Sketch instead keeps the ImageNet label space but replaces natural
photographs with black-and-white sketches~\cite{wang2019learning}. It represents
a stronger shift than ImageNet-V2 because texture, color, and photographic
detail are largely absent while object shape remains central.

\begin{figure}[t]
    \centering
    \setlength{\tabcolsep}{2pt}
    \renewcommand{\arraystretch}{1.1}
    \begin{tabular}{>{\raggedright\arraybackslash}p{0.13\textwidth} c c c}
        & \makebox[0.235\textwidth][c]{\textbf{ImageNet validation}}
        & \makebox[0.235\textwidth][c]{\textbf{ImageNet-V2}}
        & \makebox[0.235\textwidth][c]{\textbf{ImageNet-Sketch}} \\
        golden retriever
        & \makebox[0.235\textwidth][c]{\includegraphics[width=0.235\textwidth,height=0.15\textheight,keepaspectratio]{dataset_examples/imagenet_val_golden_retriever.jpeg}}
        & \makebox[0.235\textwidth][c]{\includegraphics[width=0.235\textwidth,height=0.15\textheight,keepaspectratio]{dataset_examples/imagenet_v2_golden_retriever.jpeg}}
        & \makebox[0.235\textwidth][c]{\includegraphics[width=0.235\textwidth,height=0.15\textheight,keepaspectratio]{dataset_examples/imagenet_sketch_golden_retriever.jpeg}} \\
        macaw
        & \makebox[0.235\textwidth][c]{\includegraphics[width=0.235\textwidth,height=0.15\textheight,keepaspectratio]{dataset_examples/imagenet_val_macaw.jpeg}}
        & \makebox[0.235\textwidth][c]{\includegraphics[width=0.235\textwidth,height=0.15\textheight,keepaspectratio]{dataset_examples/imagenet_v2_macaw.jpeg}}
        & \makebox[0.235\textwidth][c]{\includegraphics[width=0.235\textwidth,height=0.15\textheight,keepaspectratio]{dataset_examples/imagenet_sketch_macaw.jpeg}} \\
        strawberry
        & \makebox[0.235\textwidth][c]{\includegraphics[width=0.235\textwidth,height=0.15\textheight,keepaspectratio]{dataset_examples/imagenet_val_strawberry.jpeg}}
        & \makebox[0.235\textwidth][c]{\includegraphics[width=0.235\textwidth,height=0.15\textheight,keepaspectratio]{dataset_examples/imagenet_v2_strawberry.jpeg}}
        & \makebox[0.235\textwidth][c]{
\includegraphics[width=0.235\textwidth,height=0.15\textheight,keepaspectratio]{dataset_examples/imagenet_sketch_strawberry.jpeg}} \\
    \end{tabular}
    \caption{Examples from the ImageNet-aligned shift ladder. Each row keeps the same class while changing the collection process or visual style.}
    \label{fig:dataset_shift_severity_ladder}
\end{figure}

\paragraph{Non-ImageNet domain checks.}
Chapter 8 also uses datasets whose visual domains are not ImageNet-like. ObjectNet contains household objects photographed under controlled variations in viewpoint, rotation, and background, and was designed to expose object-recognition brittleness outside standard ImageNet biases~\cite{barbu2019objectnet}. BloodMNIST and OrganAMNIST are part of MedMNIST v2, a lightweight benchmark suite for biomedical image classification~\cite{yang2023medmnist}: BloodMNIST contains RGB microscopy images of blood cells, while OrganAMNIST contains abdominal CT slices. RVL-CDIP contains scanned document pages labeled by document type~\cite{harley2015rvlcdip}. These domains test tokenizer behavior on microscopy, medical images, and document layouts. Table~\ref{tab:non_imagenet_datasets} summarizes the datasets, and Figure~\ref{fig:non_imagenet_domain_examples} shows representative inputs.

\begin{table}[htbp]
    \centering
    \scriptsize
    \renewcommand{\arraystretch}{1.15}
    \setlength{\tabcolsep}{3pt}
    \begin{tabularx}{\textwidth}{>{\raggedright\arraybackslash}p{0.14\textwidth} >{\raggedright\arraybackslash}X >{\raggedleft\arraybackslash}p{0.11\textwidth} >{\raggedright\arraybackslash}X}
        \textbf{Dataset} & \textbf{Content} & \textbf{Images used} & \textbf{Role in Chapter 8} \\ \hline
        ObjectNet & Controlled object photos with varied viewpoints and backgrounds & 50{,}273 & Natural-object domain with different pose and background statistics. \\
        BloodMNIST & Microscopic RGB blood-cell images, 8 classes & 17{,}092 & Biomedical microscopy, dominated by local texture and cellular morphology. \\
        OrganAMNIST & Abdominal CT slices, 11 organ classes & 57{,}187 & Grayscale medical imaging with intensity statistics far from natural images. \\
        RVL-CDIP & Scanned document pages, 16 document classes & 39{,}996 & Layout- and text-like structure rather than natural-image content. \\
    \end{tabularx}
    \caption{Non-ImageNet datasets used as domain-shift checks. The image counts report the subsets used in the tokenizer experiments; for example, RVL-CDIP contains 400{,}000 images in the full source dataset.}
    \label{tab:non_imagenet_datasets}
\end{table}

\begin{figure}[t]
    \centering
    \setlength{\tabcolsep}{4pt}
    \renewcommand{\arraystretch}{1.15}
    \begin{tabular}{c c c c}
        \textbf{ObjectNet} & \textbf{BloodMNIST} & \textbf{OrganAMNIST} & \textbf{RVL-CDIP} \\
        \makebox[0.22\textwidth][c]{\includegraphics[width=0.16\textwidth,height=0.105\textheight,keepaspectratio]{dataset_examples/ood_objectnet_kettle.png}} &
        \makebox[0.22\textwidth][c]{\includegraphics[width=0.16\textwidth,height=0.105\textheight,keepaspectratio]{dataset_examples/ood_bloodmnist_class0.png}} &
        \makebox[0.22\textwidth][c]{\includegraphics[width=0.16\textwidth,height=0.105\textheight,keepaspectratio]{dataset_examples/ood_organamnist_class0.png}} &
        \makebox[0.22\textwidth][c]{\includegraphics[width=0.16\textwidth,height=0.105\textheight,keepaspectratio]{dataset_examples/ood_rvlcdip_email.png}} \\
        kettle & class 0 & class 0 & email \\
        \makebox[0.22\textwidth][c]{\includegraphics[width=0.16\textwidth,height=0.105\textheight,keepaspectratio]{dataset_examples/ood_objectnet_headphones.png}} &
        \makebox[0.22\textwidth][c]{\includegraphics[width=0.16\textwidth,height=0.105\textheight,keepaspectratio]{dataset_examples/ood_bloodmnist_class5.png}} &
        \makebox[0.22\textwidth][c]{\includegraphics[width=0.16\textwidth,height=0.105\textheight,keepaspectratio]{dataset_examples/ood_organamnist_class5.png}} &
        \makebox[0.22\textwidth][c]{\includegraphics[width=0.16\textwidth,height=0.105\textheight,keepaspectratio]{dataset_examples/ood_rvlcdip_invoice.png}} \\
        headphones & class 5 & class 5 & invoice
    \end{tabular}
    \caption{Representative non-ImageNet domain examples. ObjectNet remains object-centric, while BloodMNIST, OrganAMNIST, and RVL-CDIP move to biomedical or document domains.}
    \label{fig:non_imagenet_domain_examples}
\end{figure}

\section{Metrics}
\label{sec:metrics}

A tokenizer can be evaluated by reconstruction quality, codebook usage, and its
response to perturbations. Token frequencies determine the distribution that
the autoregressive model must learn. Perturbation response can be measured by
how far changes spread through the token grid or decoded image.

We group the metrics by what they measure: paired-image reconstruction,
image-set distributions, codebook usage, and perturbation locality. Paired
metrics compare each reconstruction with its source image. Distribution metrics
compare the aggregate distributions of two image sets; FID and related metrics
perform this comparison in a particular learned feature space. Codebook metrics
measure which codes are used and how concentrated their frequencies are.
Perturbation metrics compare matched outputs before and after an intervention,
either globally or within specified spatial regions. Figure~\ref{fig:metric_taxonomy_matrix}
summarizes these four families.

\begin{figure}[t]
    \centering
    \begingroup
\scriptsize
\renewcommand{\arraystretch}{1.16}
\setlength{\tabcolsep}{3.5pt}
\begin{tabularx}{\textwidth}{>{\raggedright\arraybackslash}p{0.18\textwidth} >{\raggedright\arraybackslash}p{0.25\textwidth} >{\raggedright\arraybackslash}p{0.25\textwidth} >{\centering\arraybackslash}X}
\rowcolor{blue!9}
\textbf{Metric family} & \textbf{Quantities used} & \textbf{Comparison unit} & \textbf{Interpretation} \\ \hline
Paired-image reconstruction & PSNR, SSIM, LPIPS & One input $x$ and its reconstruction $\hat{x}$ & PSNR, SSIM $\uparrow$; LPIPS $\downarrow$ \\
Image-set distributions & FID, sFID, precision, recall, IS & Reference and reconstruction sets & distribution agreement; IS is diagnostic \\
Codebook usage & Active codes, entropy, perplexity, top-$m$ mass & All token assignments in an evaluation set & descriptive, not a quality score \\
Perturbation locality & Flip rate, inside/outside response, leakage ratio & Matched outputs and spatial regions & less outside response means greater locality \\
\end{tabularx}
\endgroup

    \caption{Metric taxonomy used in this thesis. Evaluation combines sample-level reconstruction metrics, distribution-level image metrics, codebook-usage measures, and locality-oriented measures.}
    \label{fig:metric_taxonomy_matrix}
\end{figure}

\subsection{Sample-Level Reconstruction Metrics}
\label{subsec:sample_reconstruction_metrics}

Full-reference metrics compare each reconstruction with its corresponding
original image. Mean squared error (MSE) and peak signal-to-noise ratio (PSNR)
are full-reference pixel metrics. For an image with $N$ scalar pixel values,
\begin{align}
    \mathrm{MSE}(x,\hat{x}) &= \frac{1}{N}\sum_{i=1}^{N}(x_i-\hat{x}_i)^2, \\
    \mathrm{PSNR}(x,\hat{x}) &= 10\log_{10}\!\left(\frac{\mathrm{MAX}^2}{\mathrm{MSE}(x,\hat{x})}\right).
\end{align}
For normalized images in $[0,1]$, $\mathrm{MAX}=1$. Lower MSE and higher PSNR indicate better pixel agreement. PSNR is widely used in image restoration but is known to be limited as a perceptual quality measure~\cite{huynh2008scope}.

SSIM compares local luminance, contrast, and structure rather than only pointwise pixel error~\cite{wang2004ssim}. In simplified form,
\begin{equation}
    \mathrm{SSIM}(x,y) = l(x,y)c(x,y)s(x,y),
\end{equation}
where $l$, $c$, and $s$ denote luminance, contrast, and structure comparisons computed over local windows. SSIM values are then averaged across the image. Higher SSIM indicates stronger structural similarity.

LPIPS compares normalized deep-network activations between two images~\cite{zhang2018unreasonable}:
\begin{equation}
    \mathrm{LPIPS}(x,y)
    = \sum_l \frac{1}{H_lW_l}\sum_{h,w}
    \left\lVert
    w_l \odot \left(\hat\phi_l(x)_{h,w}-\hat\phi_l(y)_{h,w}\right)
    \right\rVert_2^2,
\end{equation}
where $\hat\phi_l$ denotes channel-normalized features at layer $l$, $w_l$
contains learned channel weights, and the distance is averaged over the
$H_l\times W_l$ spatial locations. Lower LPIPS indicates greater perceptual
similarity. Unlike MSE and PSNR, LPIPS is calibrated against human perceptual
judgments. The reconstruction evaluator used here uses its VGG-16 backbone.

\subsection{Distribution-Level Image Metrics}
\label{subsec:distribution_metrics}

FID compares two image sets by embedding them with an Inception network and fitting a Gaussian to each feature distribution~\cite{heusel2017gans}. For reference features with mean and covariance $(\mu_r,\Sigma_r)$ and reconstruction features with $(\mu_g,\Sigma_g)$,
\begin{equation}
    \mathrm{FID}
    =
    \lVert \mu_r-\mu_g\rVert_2^2
    +
    \mathrm{Tr}\!\left(\Sigma_r+\Sigma_g-2(\Sigma_r\Sigma_g)^{1/2}\right).
\end{equation}
Lower FID indicates closer feature distributions. In the literature on VQ
tokenizers, VQGAN and LlamaGen report reconstruction FID or rFID to compare the
distribution of reconstructions with original images
~\cite{esser2021taming,sun2024llamagen}.

sFID is a spatial variant used in common diffusion-model evaluator pipelines~\cite{dhariwal2021diffusion,openaiguideddiffusion,openaiguideddiffusionevals}. It uses spatial features rather than only pooled Inception features, so it can be more sensitive to spatial layout differences. It carries the same finite-sample and feature-extractor caveats as FID.

Inception Score uses an Inception classifier~\cite{salimans2016improved}:
\begin{equation}
    \mathrm{IS}
    =
    \exp\!\left(
        \mathbb{E}_{x\sim p_g}
        \left[
            \mathrm{KL}\!\left(p(y\mid x)\,\|\,p(y)\right)
        \right]
    \right).
\end{equation}
Here, $p(y\mid x)$ is the classifier's conditional class distribution for one
reconstruction and $p(y)$ is the marginal class distribution across the set.
The score increases when individual predictions are confident and predictions
across the set cover different classes. It does not measure whether a
reconstruction matches its input, so we report it as a classifier-based
diagnostic~\cite{barratt2018note,chong2019unbiased}.

For reconstructed image sets, precision measures how much of the reconstruction
distribution lies near the reference feature distribution, while recall
measures how much of the reference distribution is covered
~\cite{kynkaanniemi2019improved}. Figure~\ref{fig:precision_recall_tradeoff}
illustrates this distinction.

\begin{figure}[t]
    \centering
    \includegraphics[width=0.9\textwidth]{ch03_fig_3_3_precision_recall_tradeoff.png}
    \caption{Fidelity and coverage manifolds for Precision and Recall. Precision measures the portion of generated support within the real data manifold; Recall measures coverage of the real manifold. Source: reproduced from Kynk\"a\"anniemi et al.~\cite{kynkaanniemi2019improved}.}
    \label{fig:precision_recall_tradeoff}
\end{figure}

\subsection{Codebook Usage Metrics}
\label{subsec:codebook_usage_metrics}

Let $n_k$ be the number of assignments to code $k$ over an evaluation set, let $N_t=\sum_{k=1}^{K}n_k$ be the total number of token assignments, and let
\begin{equation}
    p_k = \frac{n_k}{N_t}.
\end{equation}
$p_k$ is the fraction of all observed token assignments that use code $k$.
The raw active-code count is
\begin{equation}
    A_{\mathrm{raw}} = \sum_{k=1}^{K}\mathbf{1}[n_k>0],
\end{equation}
and the active fraction is
\begin{equation}
    a_{\mathrm{raw}} = \frac{A_{\mathrm{raw}}}{K}.
\end{equation}
The active-code count is the number of distinct codes used at least once;
dividing it by $K$ gives the active fraction. Because a single assignment makes
a code active under this definition, Chapter~\ref{ch:rq1} also reports two
support measures. The effective active set counts codes with
$p_k\geq10^{-6}$, which is at least 13 of the 12.8 million assignments in that
evaluation. The cumulative support $C_{99}$ is the smallest frequency-ranked
set of codes covering 99\% of assignments. These measures are distinct from
perplexity, which summarizes the full frequency distribution.

Entropy and perplexity summarize how spread out the token distribution is:
\begin{align}
    H(p) &= -\sum_{k=1}^{K} p_k\log p_k, \\
    P(p) &= \exp(H(p)).
\end{align}
The intuitive meaning of perplexity is effective vocabulary size. If the tokenizer used exactly $M$ codes uniformly, then $P(p)=M$. A tokenizer may therefore have $K=16{,}384$ nominal entries but a much smaller effective vocabulary if the distribution is concentrated.

Top-$m$ mass measures concentration directly:
\begin{equation}
    M_m = \sum_{k\in \mathrm{Top}_m(p)} p_k,
\end{equation}
where $\mathrm{Top}_m(p)$ is the set of the $m$ most frequent codes. High top-$m$ mass means that a small set of codes accounts for a large fraction of assignments.

Finally, positional entropy computes the same entropy quantity at each grid position $(u,v)$ across the dataset. It asks whether different spatial positions use similar parts of the codebook or whether some positions specialize to particular tokens.

\subsection{Locality and Robustness Metrics}
\label{subsec:locality_robustness_metrics}

Perturbation metrics compare tokenizer outputs before and after a controlled
change. A global response can be measured as the change in reconstruction or
codebook statistics. Spatial locality uses a mask to divide the response into
inside and outside regions. For token position $i$, define a flip indicator
\begin{equation}
    F_i = \mathbf{1}[t_i \ne t'_i],
\end{equation}
where $t_i$ and $t'_i$ are the clean and perturbed token assignments. For a set of token positions $R$,
\begin{equation}
    r(R) = \mathbb{E}_{i\in R}[F_i].
\end{equation}
For a local image perturbation, the leakage ratio compares outside-region and
inside-region flip rates:
\begin{equation}
    \lambda_{\mathrm{enc}} =
    \frac{r(R_{\mathrm{out}})}
         {r(R_{\mathrm{in}})+\varepsilon}.
\end{equation}
The small constant $\varepsilon$ avoids division by zero when no inside-patch flips occur. A low leakage ratio means that token changes are mostly confined to the perturbed region; a high value means that the encoder or quantizer has propagated the perturbation to distant token positions.

The same definition applies to decoder locality after replacing the token flip
indicator with an image-change measure. Later chapters specify the perturbation,
mask, and response measure used by each experiment.

\begin{figure}[t]
    \centering
    \resizebox{0.9\textwidth}{!}{\begin{tikzpicture}[
    font=\sffamily\scriptsize,
    arrow/.style={-{Latex[length=2mm]}, semithick, black!65},
    eq/.style={draw=black!35, rounded corners=2pt, fill=black!2,
               text width=57mm, align=left, inner sep=6pt}
]
\begin{scope}[xshift=-58mm,yshift=-19mm]
    \fill[blue!7] (0,0) rectangle (44mm,44mm);
    \fill[red!18] (15mm,15mm) rectangle (29mm,29mm);
    \draw[step=7.333mm, black!45] (0,0) grid (44mm,44mm);
    \draw[red!70!black, very thick] (15mm,15mm) rectangle (29mm,29mm);
    \node[font=\bfseries, above] at (22mm,44mm) {spatial response grid};
    \node[red!70!black, align=center] at (22mm,22mm) {$R_{\mathrm{in}}$\\intervention};
    \node[blue!60!black, align=center] at (7mm,37mm) {$R_{\mathrm{out}}$};
\end{scope}

\node[eq] (flip) at (29mm,15mm)
    {\textbf{Flip rate in a region}\\[1mm]
     $F_i=\mathbb{1}[t_i\neq t_i']$,\qquad
     $r(R)=\dfrac{1}{|R|}\sum_{i\in R}F_i$};
\node[eq, below=5mm of flip] (leak)
    {\textbf{Encoder leakage ratio}\\[1mm]
     $\lambda_{\mathrm{enc}}=
     \dfrac{r(R_{\mathrm{out}})}{r(R_{\mathrm{in}})+\varepsilon}$};
\node[eq, below=5mm of leak] (interpret)
    {\textbf{Interpretation}\\Low outside response means that the effect remains local.
     Decoder locality uses the analogous inside/outside partition in image space.};
\draw[arrow] (-27mm,3mm) -- (flip.west);
\end{tikzpicture}
}
    \caption{Locality and leakage as measured in this thesis. Encoder locality compares clean and perturbed token assignments inside and outside the image patch mapped to token space. Decoder locality uses the analogous partition of the decoded image response.}
    \label{fig:locality_leakage_curve}
\end{figure}

\subsection{Metric Comparison Summary}
\label{subsec:metric_comparison_summary}

Table~\ref{tab:metric_comparison_summary} summarizes the metrics used in the
thesis. Paired metrics compare images, distribution metrics compare image sets,
codebook metrics describe token frequencies, and locality metrics describe the
response to controlled perturbations.

\begin{table}[t]
    \centering
    \scriptsize
    \renewcommand{\arraystretch}{1.18}
    \setlength{\tabcolsep}{3pt}
    \begin{tabularx}{\textwidth}{>{\raggedright\arraybackslash}p{0.12\textwidth} >{\raggedright\arraybackslash}p{0.20\textwidth} >{\raggedright\arraybackslash}p{0.18\textwidth} >{\raggedright\arraybackslash}p{0.17\textwidth} >{\raggedright\arraybackslash}X}
        \textbf{Metric} & \textbf{What it measures} & \textbf{Level} & \textbf{Direction} & \textbf{Main caveat} \\ \hline
        MSE / PSNR & Aligned pixel error & Sample image pair & MSE $\downarrow$, PSNR $\uparrow$ & Poor proxy for perceptual similarity. \\
        SSIM & Local luminance, contrast, and structure & Sample image pair & $\uparrow$ & Sensitive to alignment and window choices. \\
        LPIPS & Deep-feature perceptual distance & Sample image pair & $\downarrow$ & Depends on backbone and calibration. \\
        FID / sFID & Feature-distribution match between image sets & Image set & $\downarrow$ & Feature extractor, finite-sample, and Gaussian-assumption effects. \\
        IS & Classifier confidence and predicted-class diversity & Image set & $\uparrow$ & Does not directly compare outputs to real images. \\
        Precision / Recall & Fidelity to and coverage of a reference feature manifold & Image set & both $\uparrow$ & Depends on feature space and neighborhood construction. \\
        Active count / fraction & How much of the nominal codebook appears & Token distribution & $\uparrow$ & Raw activity is sensitive to rare one-off codes. \\
        Entropy / Perplexity & Effective spread of code usage & Token distribution & $\uparrow$ & Can hide spatial specialization. \\
        Top-$m$ mass & Concentration in frequent codes & Token distribution & $\downarrow$ & Choice of $m$ must be reported. \\
        Flip rate / leakage ratio & Token or image changes under controlled perturbation & Perturbation response & $\downarrow$ & Depends on perturbation mask and region definition. \\
    \end{tabularx}
    \caption{Compact comparison of metric families used in the thesis. Arrows indicate the preferred direction under a fixed evaluation protocol.}
    \label{tab:metric_comparison_summary}
\end{table}
